% SPDX-License-Identifier: Apache-2.0
% Plain article, set in XCharter to exercise `texlive_archives`.
%
% XCharter is in no TinyTeX variant, so it can only come from an added package.
% It is also the interesting case for the mechanism: it ships a font map, which
% has to be registered with updmap-sys, and it needs three support packages
% that `texlive_archives` will not resolve on its own.
\documentclass{article}

\usepackage[T1]{fontenc}
% lmodern first, for a Type 1 typewriter face: T1 \texttt with neither lmodern
% nor cm-super present makes the engine generate a bitmap with METAFONT, which
% lands in the PDF as a Type 3 font. XCharter then overrides the roman only.
\usepackage{lmodern}
\usepackage{XCharter}
\usepackage[hidelinks]{hyperref}

\begin{document}

\section*{A Short Note}

Two documents, one toolchain. This one is set in XCharter, which reaches the
vendored TeX tree through \texttt{texlive\_archives}: four archives pinned by
sha256 in \texttt{MODULE.bazel}, unpacked into \texttt{texmf-dist}, with
\texttt{mktexlsr} and \texttt{updmap-sys} run afterwards so the engine can find
the files and embed the glyphs.

That the glyphs below are XCharter and not the default roman is the test. A
missing map would still typeset, and the letters would come out wrong.

Both PDFs are concatenated into \texttt{collection.pdf}, which is where the
other three vendored tools come in: page counts, the merge, and the outline.

\end{document}
